% 2. Hafta
% -> Scalar Objects - Variables (int, float, bool, NoneType, str)
% -> Printing, numeric operators, bool operators, operator precedence, type conversion
% -> if-else, elif

\input{../KU-Beamer-Template/style/template.tex}
\usepackage{../KU-Beamer-Template/style/koc} 
\usepackage{minted}
\usepackage{upquote}
 
\title{KOLT Python} 
\subtitle{Basic Operators \& Branching} 
\newdate{date}{25}{02}{2019}
\date{\displaydate{date}}
\author{Ahmet Uysal}

\titlegraphic{\includegraphics[scale=0.2]{../KU-Beamer-Template/style/images/logo_kolt.eps}}

\setbeamercovered{invisible} % transparent

\begin{document}
    \maketitle
    \frame{\frametitle{Agenda}\tableofcontents}

    \section{Recap}

        \begin{frame}{Comments}
            \LARGE
            \inputminted[frame=single,framesep=2pt]{python3}{../Lecture1/code-examples/comments.py}
            Python will basically ignore comments, they are purely written \textbf{for humans}!
        \end{frame}

        \begin{frame}{Variables}
            \LARGE
            \begin{table}[]
                \resizebox{0.8\textwidth}{!}{
                \begin{tabular}{|l|l|l|}
                \hline
                \textcolor{koc}{\textbf{Type}} &\textcolor{koc}{\textbf{Explanation}} & \textcolor{koc}{\textbf{Examples}} \\ \hline
                \textbf{\texttt{int}}  & represent \textbf{integers} & 3, 4, 17, -10 \\ \hline 
                \textbf{\texttt{float}} & represent \textbf{real numbers} & 3.0, 1.11, -109.123123 \\ \hline
                \textbf{\texttt{bool}} & represent \textbf{boolean} truth values & \texttt{True}, \texttt{False} \\ \hline
                \textbf{\texttt{str}} & A sequence of characters. & \textquotesingle Hello\textquotesingle, \textquotesingle \textquotesingle, \textquotesingle 3\textquotesingle \\ \hline
                \textbf{\texttt{NoneType}}& special and has one value, None & \texttt{None} \\ \hline
                \end{tabular}}
            \end{table}

            \begin{itemize}
                \item How to create a variable?        
                \texttt{variable\_name = value}
                \item How about types?
                use \texttt{type()}
                \item Can a variable change type?
                \textbf{Yes!} Just assing a new value with any type.
                \item What if we if want to convert a value between types, i.e, \textquotesingle 2\textquotesingle $\to$ 2
            \end{itemize}
        \end{frame}

        \begin{frame}{Casting}
            \LARGE
            \begin{itemize}
                \item \texttt{int(\textquotesingle 2\textquotesingle)} $\to$ 2
                \item Any possible reasons for casting?
                -taking user input
                -reading numbers from a file?
                \item Can we cast every value to every type?
                \textbf{NO!} try \texttt{int(\textquotesingle hello\textquotesingle)}
            \end{itemize}
        \end{frame}

        \begin{frame}{Console I/O(Input/Output)}
            \huge
            \textbf{\texttt{print(*args, sep=\textquotesingle \ \textquotesingle, end=\textquotesingle \textbackslash n\textquotesingle )}}
            \begin{itemize}
                \LARGE
                \item Can take arbitrary number of arguments
                \item Separates elements with space by default
                \item Adds newline character \texttt{\textquotesingle \textbackslash n\textquotesingle} to end by default
            \end{itemize}
            
            \textbf{\texttt{input([prompt])}}
            \begin{itemize}
                \LARGE
                \item Prints the prompt to Console
                \item Program is paused until user enters something
                \item \textbf{returns an \texttt{str} object!} 
            \end{itemize}
        \end{frame}

    \section{Basic Operators}
        \begin{frame}{Truthy \& Falsy Values}
            \LARGE
            Last week, we tried to cast some values to \texttt{bool}, let's look at them in detail.
            \pause
            \begin{columns}
                \column{0.45\textwidth}
                    \inputminted[firstline=1, lastline=8, frame=single,framesep=2pt]{python3}{code-examples/truthy-falsy.py}
                \pause
                \column{0.54\textwidth}
                    \inputminted[firstline=9, lastline=16, frame=single,framesep=2pt]{python3}{code-examples/truthy-falsy.py}
            \end{columns}
        \end{frame}

        \begin{frame}{\texttt{bool} Operators}
            \LARGE
            How to represent logical operations in Python? (and, or, not)
            \begin{table}[]
                \resizebox{0.45\textwidth}{!}{
                \begin{tabular}{|c|c|c|c|c|}
                \hline
                \textcolor{koc}{\textbf{A}} &\textcolor{koc}{\textbf{B}} & \textcolor{koc}{\textbf{A or B}} & \textcolor{koc}{\textbf{A and B}} & \textcolor{koc}{\textbf{not A}}\\ \hline
                \textbf{\texttt{True}}  & \textbf{\texttt{True}}  & \textbf{\texttt{True}}  & \textbf{\texttt{True}} & \textbf{\texttt{False}} \\ \hline
                \textbf{\texttt{True}} & \textbf{\texttt{False}}  & \textbf{\texttt{True}} & \textbf{\texttt{False}} & \textbf{\texttt{False}} \\ \hline
                \textbf{\texttt{False}} & \textbf{\texttt{True}}  & \textbf{\texttt{True}} & \textbf{\texttt{False}} & \textbf{\texttt{True}} \\ \hline
                \textbf{\texttt{False}} & \textbf{\texttt{False}} & \textbf{\texttt{False}} & \textbf{\texttt{False}} & \textbf{\texttt{True}} \\ \hline
                \end{tabular}}
            \end{table}
            \pause
            \begin{columns}
                \column{0.25\textwidth}
                    \begin{itemize}
                        \item \textbf{\texttt{and}}
                        \pause
                        \item \textbf{\texttt{or}}
                        \pause
                        \item \textbf{\texttt{not}}
                    \end{itemize}
                \pause
                \column{0.75\textwidth}
                    \texttt{True or False and False} $\Rightarrow$ 
                    \pause
                    \textbf{\texttt{True}} \\
                    \pause
                    \Huge
                    \textbf{WHY?}
            \end{columns}
            
        \end{frame}

        \begin{frame}{Operator Precedence}
            \LARGE
            Logical operators are evaluated in this order:
            \begin{enumerate}
                \item \texttt{not}
                \item \texttt{and}
                \item \texttt{or}
            \end{enumerate}
            You can override this order with parentheses\\
            \pause
            \texttt{(True or False) and False} $\Rightarrow$ 
            \pause
            \textbf{\texttt{False}} \\
        \end{frame}

        \begin{frame}{Short-Circuit Evaluation}
            \LARGE
            \texttt{X}: Any boolean value\\
            \texttt{True or X} $\Rightarrow$ 
            \pause
            \textbf{\texttt{True}}\\
            \pause
            \texttt{False and X} $\Rightarrow$ 
            \pause
            \textbf{\texttt{False}}\\
            \pause
            Python is smart enough to take advantage of this!
            \pause
            \inputminted[frame=single,framesep=2pt]{python3}{code-examples/short_circuit.py}
        \end{frame}

        \begin{frame}{Arithmetic Operators}
            \LARGE
            These operations are applicable on Numeric types: \texttt{int} and \texttt{float}
            \pause
            \begin{columns}
               \begin{column}{0.35\textwidth}
                \vspace{-5mm}
                \begin{itemize}
                    \item \texttt{+}: Addition
                    \pause
                    \item \texttt{-}: Subtraction
                    \pause
                    \item \texttt{*}: Multiplication
                    \pause
                    \item \texttt{/}: Division
                    \pause
                    \item \texttt{//}: Floor(integer) Division
                    \pause
                    \item \texttt{\%}: Modulo
                    \pause
                    \item \texttt{**}: Power 
                \end{itemize}
               \end{column}
               \pause
               \begin{column}{0.65\textwidth}
                \inputminted[frame=single,framesep=2pt]{python3}{code-examples/numeric_operators.py}
               \end{column} 
            \end{columns}
        \end{frame}

        \begin{frame}{Comparison Operators}
            \LARGE
            \begin{columns}
                \begin{column}{0.4\textwidth}
                    \begin{itemize}
                        \item \texttt{<}: Strictly less than
                        \pause
                        \item \texttt{<=}: Less than or equal
                        \pause
                        \item \texttt{>}: Strictly greater than
                        \pause
                        \item \texttt{>=}: Greater than or equal
                        \pause
                        \item \texttt{==}: Equal
                        \pause
                        \item \texttt{!=}: Not equal
                    \end{itemize}
                \end{column}
                \pause
                \begin{column}{0.6\textwidth}
                    \inputminted[frame=single,framesep=2pt]{python3}{code-examples/comparison.py}
                \end{column}
            \end{columns}
        \end{frame}

        \begin{frame}{Chained Comparisons}
            \LARGE  
            \texttt{1 < 2 < 3} $\Rightarrow$ 
            \pause
            \textbf{\texttt{True}}\\
            \pause
            You can chain arbitrarily many comparison operations together.\\
            \pause
            $v_{i}$: variables/values, $op_{i}$: comparison operators\\
            \pause
            \texttt{$v_{1}\ op_{1}\ v_{2}\ op_{2}\ v_{3}\ ...\ op_{n-1}\ v_{n}$} is equivalent to:\\
            \texttt{$v_{1}\ op_{1}\ v_{2}$ \textbf{and} $v_{2}\ op_{2}\ v_{3}$ \textbf{and} ...$v_{n-1}\ op_{n-1}\ v_{2}$}\\
            \pause
            \vspace{-4mm}
            \inputminted[frame=single,framesep=2pt]{python3}{code-examples/chained.py}
        \end{frame}

        \begin{frame}{Assignment Operators}
            \LARGE
            We have already seen \texttt{\textquotesingle=\textquotesingle}: 
            \pause
            \texttt{variable\_name = value}


        \end{frame}

    \section{Branching}
        \begin{frame}{Branching}
            \vspace{-3mm}
            \begin{columns}
                \column{0.5\textwidth}
                \inputminted[firstline=1, lastline=4, frame=single,framesep=2pt]{python3}{code-examples/branching.py}
                \inputminted[firstline=6, lastline=13, frame=single,framesep=2pt]{python3}{code-examples/branching.py}
                \column{0.5\textwidth}
                \inputminted[firstline=15, lastline=27, frame=single,framesep=2pt]{python3}{code-examples/branching.py}
            \end{columns}
            \begin{itemize}
                \item \texttt{<condition>} has a \textbf{\texttt{bool}} value (\texttt{True} or \texttt{False})
                \item Which expressions will be evaluated in which conditions?
            \end{itemize}
        \end{frame}

\end{document}